
\begin{problem}[2026年博知汇BWPhO][通电平板中的电荷密度]{35}
    考虑空间直角坐标系 $O$-$xyz$ 中填充范围为 $-d \leq z \leq d$ 的无限大导体平板，其电导率为 $\sigma$。现在沿 $x$ 轴正方向通入电流，$y$ 方向上单位长度通入的电流大小为 $I$。已知该导体内的载流子只有电子一种，其数密度为 $n$。为了与自然常数 $e$ 作区分，以 $q$ 表示元电荷 $(q > 0)$。在本题中，始终认为导体平板内的净电荷密度 $\rho \ll nq$。以下第（2）（3）（6）问保留至最低阶非零项，第（4）（5）问保留至含 $\omega$ 的最低阶修正项。
    \begin{subproblem}
        考虑 $I = I_0$ 恒定，且体系达到稳定的情形，试证明此时平板内的电流密度 $\pmb{j}$ 与 $z$ 无关（即为均匀的）。
    \end{subproblem}
    \begin{subproblem}
        承第（1）问，由于导体内的电流在导体内部产生的磁感应强度，导体内会产生净电荷密度。试求出此时导体内的净电荷密度分布 $\rho(z)$。
    \end{subproblem}
    \begin{subproblem}
        若电流的通入是突然的，且开始通入后单位长度上电流 $I = I_{0}$ 恒定，即
        \begin{equation}
            I (t) = \left\{ \begin{array}{l l} 0, & t <   0 \\ I _ {0}, & t \geq 0 \end{array} \right.
        \end{equation}
        试求出导体内净电荷密度分布随电流通入时长 $t$ 的变化关系 $\rho (z,t)$。（计算磁场分布时忽略 $z$ 方向电流的影响，下同）
    \end{subproblem}
    \begin{subproblem}
        现在考虑交流情形，取 $I = I_0\cos \omega t$。此时导体平板上将发生趋肤效应。设 $\omega \varepsilon_0 / \sigma \ll \sqrt{\omega \mu_0 \sigma} d \ll 1$（从而可以忽略位移电流对磁场的贡献）。试求解稳定情况下导体内的电流密度的 $x$ 分量大小 $j$ 与 $z$ 的关系 $j(z)$。（这里“稳定”是指随时间周期性变化，下同）
    \end{subproblem}
    \begin{subproblem}
        承第（4）问，求解稳定情况下导体内的净电荷密度分布随时间 $t$ 的变化关系 $\rho (z,t)$。
    \end{subproblem}
    \begin{subproblem}
        承第（5）问，但不考虑趋肤效应，求解因为体电荷分布变化导致的额外焦耳热功率与只考虑 $x$ 方向电流求得的焦耳热功率之比 $\eta(t)$。
    \end{subproblem}
\end{problem}
